\section{Introduction to the Field of 3D Reconstruction}
3D Gaussian Splatting is a new technology in the field of 3D reconstruction which requires introducing certain basic terminology.

\textit{3D reconstruction} marks a field in Computer Graphics which allows creating 3D \textit{scenes} of different kinds from a given input like a set of images or frames of a video.
Once this 3D scene has been created, \textit{novel-view synthesis} refers to the process of gaining previously non-existing imagery from this 3D scene~\cite{liuSurvey3DReconstruction2025, tewariStateArtNeural2020}. However, both processes are commonly aggregated into the single term 3D reconstruction, a practice we do not adopt here.
Next, \textit{real-time novel-view synthesis} is the ability to synthesize new imagery at a rate of 24 to 30 frames per second (FPS) or more~\cite{heSurvey3DGaussian2026}, making it possible to smoothly explore a reconstructed scene comparable to commonly used non-reconstructed mesh-based scenes found e.g. in 3D gaming.

First efforts in this field date back to the 1990s with work on multi-view geometry~\cite{liuSurvey3DReconstruction2025} and remained a niche area of research due to lack of powerful models until the general emergence of neural networks in the late 2010s. 
Neural networks for the first time allowed creating 3D scenes with realistic image quality using the radiance field approach.
An \textit{implicit radiance field}~\cite{hubertEditingImplicitExplicit2025a} is based on the mathematical concept of the five-dimensional plenoptic function $L: x, y, z, \delta, \phi \rightarrow c, \alpha$ which expresses a scene by continuously mapping any given point $(x, y, z)$ in space and a set of viewing angles $\delta$ and $\phi$ (pitch and yaw) to a color $c$ and an opacity $\alpha$ at that point.
\begingroup
\clubpenalty=0
\widowpenalty=0
\textit{Explicit radiance fields} restrict the plenoptic function to a finite set of points, for example using a rigid 3D \textit{voxel} grid.
\par
\endgroup

One well-known implementation of radiance fields are \textit{Neural Radiance Fields} (NeRFs) introduced by Mildenhall et al. in 2020~\cite{mildenhallNeRFRepresentingScenes2020}. NeRFs take a set of input images and use a neural network to learn an implicit radiance field representation of the resulting scene, which means the full scene is encoded in the neural network's weights.
In order to perform novel-view synthesis for a NeRF, a technique called \textit{volumetric ray marching} is employed. 
For every pixel in the novel view the NeRF algorithm transmits a ray through the implicit radiance field resulting in a very high number of queries to the neural network for each (discretized) point along the path of the ray before adding up all predicted RGB and opacity values to form the final color value of each pixel~\cite{mildenhallNeRFRepresentingScenes2020}.
For more information on NeRFs refer to the previous paper in this seminar.

The other aspect which is important to know is the concept of \textit{splatting}. First introduced in 1991 by Westover~\cite{Westover1991SplattingAP} and refined later in the early 2000s by a group around Zwicker~\cite{zwickerEWAVolumeSplatting2001, zwickerSurfaceSplatting2001}, it takes the reverse approach to ray marching by projecting all elements in the 3D scene onto the view plane towards the camera, a process which colloquially can be called ``splatting'' or ``squashing''. 
It was also Zwicker et al. in the same publications who introduced \textit{3D Gaussians} as a data structure to simplify and extend splatting.

Lastly, there are important basic concepts in Computer Graphics that need explaining especially in the context of splatting, starting with \textit{world space}, \textit{view space}, and \textit{image space}. The first describing the space in which the 3D scene is built; the second describing the same space shifted with the camera at the origin; and the third describing the space after \textit{camera projection}.
Then there is the near plane, also called the \textit{view plane} or projection plane and the \textit{far plane} as well as the \textit{camera pose} or position. The camera pose and its \textit{field of view} (FOV) decide which elements of the scene between the view plane and the far plane are to be included in the projection onto the view plane.
This setup is known as the ``pyramid of vision'' or \textit{view frustum} as it takes the shape of a frustum, see Figure~\ref{fig:view_frustum} on page~\pageref{fig:view_frustum}. \textit{Frustum culling} removes all elements outside the view frustum. For more general aspects of computer graphics see~\cite{marschnerFundamentalsComputerGraphics2021}.